\section{The $48$-vertex floor and the degree-twelve terminal}
\label{sec:h12}

\subsection{Minimum degree in $\mathcal R(4,48)$}

\begin{lemma}
\label{lem:R448-degree}
Every graph $J\in\mathcal R(4,48)$ has minimum degree at least $12$.
\end{lemma}

\begin{proof}
A vertex of degree at most $10$ leaves at least $37$ nonneighbors, containing a forbidden member of $\mathcal R(3,37)$. Suppose instead that $p$ has degree $11$. Its $36$-vertex nonneighborhood $H$ lies in $\mathcal R(3,36)$, so \cref{lem:R336} gives $e(H)=234$. Each vertex of the $11$-vertex shell has at most three neighbors in $H$, since otherwise the corresponding $37$-set would exceed $D(37)=237$.

The shell cannot be a clique, because together with $p$ it would form a $K_{12}$. Choose nonadjacent shell vertices $a,b$. Their common antineighborhood in $H$ has independence number at most $2$, since an independent triple there together with $a,b$ would be an independent five-set. By \cref{lem:R224}, that common antineighborhood has at most $23$ vertices. Thus the union of the two rows into $H$ has at least $36-23=13$ vertices. But each row has size at most $3$, so their union has size at most $6$, a contradiction.
\end{proof}

\subsection{The charge-$20$ shell lemma}

Consider a degree-$12$ pivot in a parent graph whose relevant $48$-vertex cores all have minimum degree at least $12$. Let $M$ be the missing graph on the $12$-vertex shell. For a shell vertex $a$, write
\[
g_a=d_F(a)-12\geq0,\qquad t_a=d_M(a)+g_a,
\qquad q=e(M)+\sum_a g_a.
\]
The inherited conditions are
\begin{equation}
e(M)\geq11,\quad \alpha(M)\leq10,\quad
t_a+t_b\geq12\ \ (ab\in E(M)),\quad
t(T)\geq24\ \ (T\text{ a triangle of }M).
\label{eq:q20-conditions}
\end{equation}
The first inequality follows from the $13$-vertex cap $D(13)=67$, since the closed neighborhood has $78-e(M)$ edges. The second follows because a target $K_{11}$ in the shell together with the pivot would form a $K_{12}$. The pair and triangle row-union inequalities follow from \cref{lem:R224} and the clique cap.

\begin{lemma}[Charge-$20$ classification]
\label{lem:q20}
Under \eqref{eq:q20-conditions}, $q\geq20$. If $q=20$, then
\[
M=K_{2,r}\ \dot\cup\ (10-r)K_1,\qquad 6\leq r\leq10,
\]
with the two hubs of weight $10$, hub decorations $10-r$, and all ten low vertices regular; each low row has size $2$ or $0$.
\end{lemma}

\begin{proof}
If $M$ contains a triangle $T$, then $q\geq t(T)-e_M(T)\geq24-3=21$. Suppose $M$ is triangle-free. It has a matching of size at least two: a graph with matching number at most one is a star or a triangle plus isolated vertices; the star has an independent set of size $11$, while the triangle has only three edges, contrary to \eqref{eq:q20-conditions}.

Let $S$ be the four endpoints of a two-matching. Mantel's theorem gives $e_M(S)\leq4$, while the pair-row conditions imply $t(S)\geq24$. Therefore $q\geq t(S)-e_M(S)\geq20$.

If equality holds, then $M[S]=K_{2,2}$, $t(S)=24$, and all eight other shell vertices are independent and undecorated. A three-matching would give $q\geq36-9=27$, so none exists. Triangle-freeness implies that every exterior point attaches only to one side of the $K_{2,2}$. Attachments cannot occur on both sides: an edge joining two such exterior attachments, together with the edge between the two unused cycle vertices, would produce a three-matching. Hence all attachments meet one two-point side $X$, giving a two-hub cover and an independent ten-set $Y$.

Summing the four edge inequalities across the $K_{2,2}$ shows that every such edge has endpoint weight sum $12$. The two hub weights have a common value $a$, and the two opposite weights are $12-a$. Since each opposite-side vertex has missing degree two, $a\leq10$. As $e(M)\geq11$, at least seven exterior attachments occur. Any nonisolated exterior point is regular and has missing degree at most two; its edge inequality with a hub forces $a=10$ and its row size to be two. Thus the only equality graphs are the displayed $K_{2,r}$ forms, with $6\leq r\leq10$ and the stated decorations.
\end{proof}

\subsection{The inserted-center contradiction}

\begin{lemma}
\label{lem:inserted-center}
The charge-$20$ equality configuration in \cref{lem:q20} is impossible if every degree-$12$ repivot of the parent has a $48$-vertex core of minimum degree at least $12$.
\end{lemma}

\begin{proof}
Choose a leaf $y$ of the $K_{2,r}$ shell. Repivoting at $y$ removes its two-point core row and inserts the two old hubs. The hubs are adjacent in the target graph, since there is no hub--hub edge in the missing graph $M$. Each hub has only ten neighbors in the old core. Its degree in the new core is therefore at most $10+1=11$, even before subtracting any edges into the deleted row. This contradicts the assumed minimum degree $12$ of every actual $48$-vertex core.
\end{proof}

\subsection{The $310$-edge floor}

\begin{theorem}[The $48$-vertex floor]
\label{thm:floor48}
Every $J\in\mathcal R(4,48)$ has at least $310$ edges.
\end{theorem}

\begin{proof}
By \cref{lem:R448-degree}, $\delta(J)\geq12$. If $\delta(J)\geq13$, then $e(J)\geq312$. Otherwise choose a degree-$12$ pivot. Its $35$-vertex core $H$ belongs to $\mathcal R(3,35)$, has minimum degree at least $12$, and has at least $211$ edges by \cref{lem:core211}. The shell missing graph $M$ has at least $11$ edges by $D(13)=67$, and $\alpha(M)\leq10$.

For every missing shell edge, the union of the two core rows has size at least $12$; otherwise the uncovered core would contain a member of $\mathcal R(2,24)$. For every missing triangle, the row union has size at least $24$, since the uncovered core has independence number at most one and hence is a clique of size at most $11$. Thus \cref{lem:q20} applies and gives $q\geq20$.

If $q=20$, the equality classification and \cref{lem:inserted-center} apply: a leaf repivot produces another genuine $35$-vertex core in $\mathcal R(3,35)$ with minimum degree at least $12$, contradicting the inserted-center degree bound. Therefore $q\geq21$. The exact edge partition gives
\[
e(J)=78+q+e(H)\geq78+21+211=310.
\]
\end{proof}

\subsection{Eliminating the degree-twelve $61$-vertex residual}

\begin{proposition}
\label{prop:h12-empty}
There is no graph $F$ satisfying the $61$-vertex contract \eqref{eq:T61-contract} with $\delta(F)=12$.
\end{proposition}

\begin{proof}
A degree-$12$ pivot has a genuine $48$-vertex core $C\in\mathcal R(4,48)$, so $e(C)\geq310$ by \cref{thm:floor48}. Its shell charge satisfies $q\geq20$ by \cref{lem:q20}: the pair-row union is at least $12$, and a missing triangle has row union at least $24$, since its uncovered core is a clique and has at most $11$ vertices. The exact edge partition yields
\[
E=e(F)=78+q+e(C)\geq78+20+310=408.
\]
Thus $E<408$ is impossible. At $E=408$, equality forces $q=20$ and $e(C)=310$. Every degree-$12$ repivot produces another actual $48$-vertex core in $\mathcal R(4,48)$, whose minimum degree is at least $12$ by \cref{lem:R448-degree}. The inserted-center contradiction in \cref{lem:inserted-center} now rules out the equality case as well.
\end{proof}

Finally, if a $61$-vertex residual satisfying \eqref{eq:T61-contract} had a vertex of degree at most $11$, its nonneighborhood would contain a contextual member of $\mathcal R(4,49)$, contradicting \cref{lem:context49}. If its minimum degree were at least $14$, the handshake lemma would give $E\geq\lceil61\cdot14/2\rceil=427>408$. Thus the only possible minimum degrees are $12$ and $13$; the former has just been excluded.

